\chapter{Insertion on Two Kinova Gen3 Arms}
\label{ch:kinova}

This chapter executes the insertion stage of Chapter~\ref{ch:tetris} on two physical arms. The convex program, the scalar squeeze loop, the jam sentry and the learnt routes are the ones used in simulation; what changes is the apparatus that realizes them. Each wrist carries a six-axis force/torque sensor, so the hand wrenches that the controller closes on are measured rather than estimated from joint torques.

The runs reported here are a first hardware series, recorded on 21 and 22 October 2026. The piece is presented by hand at the grasp pose and the board is built by hand from the route file. Ten insertions over four worlds are not a campaign, and further runs are planned. The chapter therefore states what was executed and what was measured, and keeps the outcomes as the operator recorded them.

\section{The cell}

Two Kinova Gen3 7-DoF arms stand with parallel bases $600$ mm apart. A shelf stands in front of them, as in Fig.~\ref{fig:kinova-cell}. Its floor runs from $120$ to $770$ mm ahead of the bases and lies $30$ mm below the base plane. The shelf is $125$ mm deep and has no front or back panel, so each tool reaches one face of the piece from its own side. All poses are expressed in the shelf frame of the simulation: origin on the floor at the far end of column~$0$, $z$ up, the shelf plane is $x$--$z$ and the pad axis is $y$. Each arm's base pose in that frame was measured by hand.

\thesisfigure[0.62]{ch11_kinova/cell_top_view.pdf}{The cell from above, in the frame of the right base. The two tools are drawn at the grasp of the T route; the piece moves in the mid-plane of the shelf, between them.}{fig:kinova-cell}

The pieces are built from $100$ mm polystyrene cubes; a whole tetromino weighs $50$ g. The routes were trained at this size (Chapter~\ref{ch:tetris}): the five-piece kit worlds with the slot narrowed to $75\%$ and the route starting with the piece's grasp point at $800$ mm, and the narrowed worlds of the column-$2$--$5$ set.

Each tool is, from the flange outwards: an ATI Mini40 (SI-80-4) force/torque sensor with its adapters, a free coil spring and a printed tri-pad of three spheres of $15$ mm radius on a $70$ mm triangle. The front of the spheres is $100$ mm from the flange. The spring is $50$ mm long, with an axial stiffness of $2700$ N/m, and is not guided. In series with a $1000$ N/m arm it gives about $730$ N/m per hand along the squeeze axis. Its estimated torsional stiffness is about $0.25$ Nm/rad, against $20$ Nm/rad at the wrist. The spring therefore passes forces but almost no moment about the pad axis, which matters for the twist channel below. The $150$ g mounted on each flange are included in the URDF used for gravity compensation.

\begin{table}[htbp]\centering\small
\caption{The hardware configuration of the runs of 21--22 October 2026.}\label{tab:kinova-setup}
\resizebox{\textwidth}{!}{\begin{tabular}{@{}ll@{}}
\toprule
Arms & 2 $\times$ Kinova Gen3, 7 DoF, bases parallel, $600$ mm apart \\
Arm control & Cartesian impedance at the flange, joint torque at $1$ kHz \\
Stiffness, damping & $1000$ N/m, $20$ Nm/rad; $5$ N\,s/m, $1$ Nm\,s/rad \\
Wrist sensors & ATI Mini40 SI-80-4, Ethernet NetBox, $1$ kHz, means of 5 at $200$ Hz \\
Tool & sensor, free spring ($2700$ N/m), tri-pad; $100$ mm, $150$ g past the flange \\
Pieces & $100$ mm polystyrene cubes, $50$ g per tetromino \\
Command loop & both flange targets from one object frame, $100$ Hz \\
Squeeze and twist & closed on the measured hand wrenches, every $20$ ms \\
Convex program & re-solved every $100$ ms in a worker thread \\
Speed & route, turns and seat $3\times$ slower than trained \\
Safety stops & $60$ N per hand, $40$ mm tracking, jam sentry, sensor silent $50$ ms \\
\bottomrule
\end{tabular}}
\end{table}

\subsection{Control on the arms}

Each arm runs the laboratory's Cartesian impedance controller: a target flange pose in, joint torques out at $1$ kHz, with gravity compensated from the URDF. The stiffness is diagonal in the base axes and fixed at launch. The values in Table~\ref{tab:kinova-setup} are those used in the bimanual friction-margin experiments on the same arms~\cite{frictioncone2026}, not the $1500$ N/m and $40$ Nm/rad of the simulated wrist.

The insertion node runs one loop at $100$ Hz. It reads both flange poses by forward kinematics and both sensor readings, and sends both flange targets. As in simulation, the targets follow from one commanded object frame plus a scalar squeeze along the pad axis. Every $20$ ms the squeeze is integrated on the measured normal force towards the convex program's target. With the full wrench enabled, the opposite twist of the two flanges about the pad axis is integrated on the measured moment error as well. The convex program is re-solved every $100$ ms. On the robot computer it takes a median $47$--$58$ ms, so it runs in a worker thread: the loop keeps streaming targets with the last solution until the next one arrives.

\subsection{The wrist sensors}

Each NetBox streams at $1$ kHz; a light reader publishes means of five samples at $200$ Hz, and the loop averages the samples of its own period. The sensor's pose follows from the flange pose and a fixed rotation. That rotation was found by three holds by hand per arm: nothing touching, a press on the pad face towards the flange, and a press down on the front of the tri-pad. The two measured force directions fix the sensor axes against the flange axes, and the moment of the downward press fixes the sign. The two sensors turned out to be rotated by about $9^\circ$ and $12^\circ$ about the flange axis. At each tare, with the tools in free air, the bias is the mean of $0.5$ s of samples less the tool's weight at the current attitude, so the weight is removed again at every attitude afterwards. The standard deviation at rest was below $0.03$ N per force axis in every run.

\section{What the arms were asked to do}

Each run follows the sequence of Chapter~\ref{ch:tetris}. The piece is presented by hand between the open tools at the route's grasp pose, above the board. The tools close until contact and squeeze, and the operator lets go. The pair then carries the piece to the mouth of the slot and along the learnt route. A slow straight seat follows, then a descent until the shelf is felt to carry the piece, the release, and a retreat. The commanded path is the one exported from the full simulation, scaled to the real cell, at one third of the trained speed.

\thesisfigure{ch11_kinova/planned_vs_executed.pdf}{The four worlds executed, as built: neighbours, the piece where it is presented (solid) and its slot (dashed). The dashed black line is the route's planned path of the grasp centre, the midpoint of the two pads. The coloured lines are the measured midpoint of the two flanges in every run of that world, from the mouth of the slot to the seat.}{fig:kinova-routes}

Figure~\ref{fig:kinova-routes} shows the four worlds. The T is presented turned by $8^\circ$ from its slot attitude, the L and the Z by $3^\circ$, the S upright; each route turns the piece into its slot attitude on the way down.

\section{What they did}

Table~\ref{tab:kinova-runs} lists the ten runs, all with the full wrench and the wrist sensors as the source of the hand wrenches. Every run followed its route, found the seat and released the piece, with no jam stop. The operator recorded eight insertions and two entries; in the first T run the neighbouring L was not pushed aside, in the second all neighbours were. The placement error, the worst cell against its target as the simulation scores it, has not yet been measured by hand.

\begin{table}[htbp]\centering\footnotesize
\caption{The runs of 21--22 October 2026. Asked and measured: the median normal force per hand while carrying (target of the convex program and wrist sensor). Side: the largest horizontal force the environment exerted on the piece. Twist: the largest rotation the twist channel added to a flange. Solve: median time of one convex program on the robot computer.}\label{tab:kinova-runs}
\setlength{\tabcolsep}{4pt}
\resizebox{\textwidth}{!}{\input{figures/ch11_kinova/runs_table.tex}}
\end{table}

\paragraph{Tracking.} The measured midpoint of the flanges stays within a median $5$--$8$ mm of its commanded position, at most $14$--$21$ mm, while carrying. This is the lag of a $1000$ N/m impedance under the contact forces and the residual of the gravity model. Close to the route start, a $0.77$ m reach with the tool horizontal, inverse kinematics on the controller's URDF finds the flange up to about $9$ mm short of the commanded pose of the T route: the start pose is at the edge of the arms' workspace, and the impedance absorbs the difference as tracking error.

\paragraph{Squeeze.} The squeeze closed on the measured force follows the convex program. The median normal force per hand differs from its target by at most $0.1$ N in every run (Table~\ref{tab:kinova-runs}). The targets are $3.9$--$4.9$ N per hand for a $0.49$ N piece, and the largest measured value is $8.0$ N. Figure~\ref{fig:kinova-wrench} shows one run of the L. During the route the target varies with the predicted slot load and the measured force follows it, with oscillations of about $1$ N at roughly $1$ Hz. These are consistent with the soft spring in series with the arm.

\thesisfigure{ch11_kinova/wrench_one_run.pdf}{One run of the L (22 October, 15:00), from contact to release; the time the arms held still waiting for the operator is cut out. Background: approach and squeeze (grey), carry to the mouth (blue), route (sand), slow seat and seat by feel (green), release (grey). From the top: normal force of each hand measured by its wrist sensor against the target of the convex program; the force the environment exerts on the piece, with the jam limit (off the scale during the route, about $40$--$60$ N); the moment each hand applies about the pad axis; the twist correction added to each flange.}{fig:kinova-wrench}

\paragraph{Contact with the board.} The horizontal force the environment exerted on the piece stayed below $5$ N in every run, far under the jam limit; the S pushed hardest. In Fig.~\ref{fig:kinova-wrench} the vertical force rises to $4.5$ N at the seat: this is the shelf taking the piece, the event the seat-by-feel step waits for.

\paragraph{Twist.} The twist channel was active in every run, but it had little to act on. The moment each hand applied about the pad axis stayed within $\pm0.04$ Nm. Meanwhile the correction rotated the flanges steadily, up to $1.4$--$3.6^\circ$, without reducing that moment. This is what the spring predicts: at about $0.25$ Nm/rad, $2.5^\circ$ of flange rotation changes the moment at the pads by about $0.01$ Nm. With this tool the opposite twist cannot be regulated through the flange. The insertions succeeded with the squeeze carrying the piece, and the twist channel contributed at most a small rotation. A tool that transmits torsion, or a twist actuated at the pad, would be needed before the channel can be evaluated.

\paragraph{The convex program.} Out of $203$--$309$ solves per run, at most one failed; that tick fell back to the trained pickup squeeze. A median solve of $47$--$58$ ms within a $100$ ms period was enough once the program left the streaming loop. Solved inline, the same program had stalled the loop for $40$--$260$ ms about once per second during carrying, and the arms visibly jerked.

\paragraph{Why the wrist sensors.} Earlier runs used the external wrench estimated from the joint torque sensors. With the arms empty, that estimate drifted by up to $8$--$9$ N per arm along a route, repeatably. Each route needed an empty-handed recording of the error, subtracted at the nearest flange position, and the moments could not be corrected that way. Figure~\ref{fig:kinova-ati-est} compares the two sources on the same run. During the route, the estimate departs from the sensor by up to about $4$ N. After the release, it settles at $-1.5$ to $-2.5$ N where the sensor reads zero. These errors are of the size of the squeeze itself.

\thesisfigure{ch11_kinova/ati_vs_estimate.pdf}{Normal force of each hand in the run of Fig.~\ref{fig:kinova-wrench}: the wrist sensor, which the controller used, against the joint-torque estimate logged beside it (tared at the same instant).}{fig:kinova-ati-est}

\section{Earlier runs with the joint-torque estimate}

Before the wrist sensors were mounted, the same insertion ran with the joint-torque estimate and the empty-handed correction. On 16 October the T was inserted with the scalar squeeze; the Z entered, but the stack of neighbours half fell. That run used the correction recorded for the T, not for the Z. On 19 October the L was inserted. On 20 October the T of the column-$2$--$5$ set was the first insertion with the twist channel enabled. It was inserted, but not smoothly. Its log shows the two hands opposing each other vertically instead of sharing the weight, with the measured forces needing more friction than the model's $\mu=1$ in two thirds of the ticks. These runs established the sequence on the arms and showed the limits of the estimate that the wrist sensors removed.

\section{What these runs show and what they do not}

They show that the routes learnt in simulation, scaled to the laboratory pieces, can be executed by two impedance-controlled arms following one object frame. They also show that a squeeze closed on wrist-sensor forces can track the convex program's target to a tenth of a newton, on a soft tool, while the piece is pushed into a narrowed slot.

They do not yet show placement accuracy, which is still to be measured, nor the effect of the full wrench: the twist channel had almost no authority through a free spring, and no run with the scalar squeeze alone was made with the wrist sensors. Each world was run two or three times with the piece presented by hand, so the variability of the grasp is the operator's. Neighbours were pushed in some runs and not in others, and how far they moved was not measured. Some routes start at the edge of the arms' reach. A version of the sequence in which the arms pick the piece themselves from a fixed place on the empty shelf has been implemented and checked offline, but not run.

\todo{Measure the placement error of every run by hand; add the remaining worlds and repeat runs; run the scalar squeeze with the wrist sensors on the same worlds for a paired comparison with the full wrench; replace the free spring with a tool that transmits torsion before evaluating the twist channel; record neighbour displacement.}
